
%%%%%%%%%%%%%%%%%%%%%%%%%%%%%%%%%%%%%%%%
%           main body
%%%%%%%%%%%%%%%%%%%%%%%%%%%%%%%%%%%%%%%%


%-------------------------------------
%	TITLE PAGE
%-------------------------------------


\title[Calculus]{\LARGE \bfseries 第一部分\quad 微积分} % The short title appears at the bottom of every slide, the full title is only on the title page
\author[\tiny Back Propagation] % Your institution as it will appear on the bottom of every slide, may be shorthand to save space
{\Large \bfseries \S 9. 反向传播算法初步
}
% \author{Singular Value Decomposition} % Your name
% \date{\today} % Date, can be changed to a custom date
\date{}

\begin{document}

\begin{frame}
\titlepage % Print the title page as the first slide
\end{frame}

% \begin{frame}
% \frametitle{内容提要} % Table of contents slide, comment this block out to remove it
% \tableofcontents % Throughout your presentation, if you choose to use \section{} and \subsection{} commands, these will automatically be printed on this slide as an overview of your presentation
% \end{frame}

%---------------------------------------
%	PRESENTATION SLIDES
%---------------------------------------

%------------------------------------------------

\begin{frame}

\tikzset{%
  every neuron/.style={
    circle,
    draw,
    minimum size=0.5cm
  },
  neuron missing/.style={
    draw=none,
    scale=2,
    text height=0.333cm,
    execute at begin node=\color{zkblue}$\vdots$
  },
}

\begin{block}{神经网络示意图}
\end{block}


\begin{center}
\begin{tikzpicture}[x=1.3cm, y=0.9cm, >=stealth]

% \foreach \m/\l [count=\y] in {1,2,3,missing,4}
%   \node [every neuron/.try, neuron \m/.try] (input-\m) at (0,2.5-\y) {$x_{\m}$};
\foreach \m/\l [count=\y] in {1,2,3}
  \node [every neuron/.try, neuron \m/.try] (input-\m) at (0,2.5-\y) {$x_{\m}$};
  \node [neuron missing/.try] (input-missing) at (0,2.5-4) {};
  \node [every neuron/.try] (input-4) at (0,2.5-5) {$x_n$};

% \foreach \m [count=\y] in {1,missing,2}
%   \node [every neuron/.try, neuron \m/.try ] (hidden-\m) at (2,2-\y*1.25) {};
  \node [every neuron/.try] (hidden-1) at (2,2-1*1.25) {$t_1$};
  \node [neuron missing/.try] (hidden-missing) at (2,2-2*1.25) {};
  \node [every neuron/.try] (hidden-2) at (2,2-3*1.25) {$t_p$};

% \foreach \m [count=\y] in {1,missing,2}
%   \node [every neuron/.try, neuron \m/.try ] (output-\m) at (4,1.5-\y) {};
  \node [every neuron/.try] (output-1) at (4,1.5-1) {$y_1$};
  \node [neuron missing/.try] (output-missing) at (4,1.5-2) {};
  \node [every neuron/.try] (output-2) at (4,1.5-3) {$y_m$};

\foreach \l [count=\i] in {1,2,3,n}
  \draw [<-] (input-\i) -- ++(-1,0)
    node [above, midway] {$I_\l$};

% \foreach \l [count=\i] in {1,n}
%   \node [above] at (hidden-\i.north) {$H_\l$};
\node [above] at (hidden-1.north) {$H_1$};
\node [below] at (hidden-2.south) {$H_p$};

\foreach \l [count=\i] in {1,m}
  \draw [->] (output-\i) -- ++(1,0)
    node [above, midway] {$O_\l$};

\foreach \i in {2,...,3}
  \foreach \j in {1,...,2}
    \draw [->] (input-\i) -- (hidden-\j);

\draw [->] (input-1) -- (hidden-2);
\draw [->] (input-4) -- (hidden-1);
\draw [->] (input-1) -- (hidden-1)
  node [above, midway] {$\color{red} w_{11}$};
\draw [->] (input-4) -- (hidden-2)
  node [below, midway] {$\color{red} w_{np}$};

% \foreach \i in {1,...,2}
%   \foreach \j in {1,...,2}
%     \draw [->] (hidden-\i) -- (output-\j);

\draw [->] (hidden-1) -- (output-2);
\draw [->] (hidden-2) -- (output-1);
\draw [->] (hidden-1) -- (output-1)
  node [above, midway] {$\color{red} v_{11}$};
\draw [->] (hidden-2) -- (output-2)
  node [below, midway] {$\color{red} v_{pm}$};

\foreach \l [count=\x from 0] in {输入层, 隐藏层, 输出层}
  \node [align=center, above] at (\x*2,2) {\l};

\end{tikzpicture}
\end{center}

\end{frame}

%------------------------------------------------

\begin{frame}

\begin{block}{神经网络}
对于以上神经网络，有
\begin{align*}
\small
t_j = \sigma( \underbrace{\sum\limits_{i=1}^n {\color{red} w_{ij}}x_i + {\color{red} b_j}}_{r_j^{(1)}} ), \quad y_j = \sigma ( \underbrace{\sum\limits_{i=1}^p {\color{red} v_{ij}}t_i + {\color{red} b'_j}}_{r_j^{(2)}})
\end{align*}
其中的$\sigma$为所谓的Activation Function, 一般是以下的Sigmoid Functions之一
\begin{itemize}
\item 逻辑函数$f(x) = \dfrac{1}{1 + e^{-x}}$,
\item 双曲函数$\tanh(x) = \dfrac{e^x - e^{-x}}{e^x + e^{-x}}$,
\item 误差函数$\operatorname{erf}(x) = \dfrac{1}{\sqrt{\pi}} \int\limits_{0}^x e^{-t^2}dt$, \; $\cdots\cdots$
\end{itemize}
\end{block}

\end{frame}

%------------------------------------------------

\begin{frame}

\begin{block}{反向传播算法原理}
预先取定一个损失函数$L$, 例如可以取
$$L(y, \widehat{y}) = (y-\widehat{y})^2$$
其中$\widehat{y}$为模型的输出值。

\vspace{0.5em}

目标：使$\mathcal{L} = \sum L(y_j, \widehat{y}_{j}) {\color{gray!50} + \text{正则化项}}$最小化。

\vspace{0.5em}

方法：梯度下降。为了记号方便，考虑没有$b,b'$的情形。
$$\mathcal{L}(v_{ij}+h_{ij}) \approx \mathcal{L}(v_{ij}) + \nabla\mathcal{L}(v_{ij})\cdot (h_{ij})$$
需要取$(v_{ij})$的增量为
$$(h_{ij}) = -\eta \nabla\mathcal{L}(v_{ij})$$
$\eta$为所谓的学习率。
\end{block}

\end{frame}

%------------------------------------------------

\begin{frame}

\begin{block}{反向传播算法原理}
用坐标表示（再加上求偏导的链式法则）就是
$$h_{ij} = -\eta\dfrac{\partial\mathcal{L}}{\partial v_{ij}} = \pause -\eta \underbrace{\dfrac{\partial L}{\partial \widehat{y}_j}}_{\partial_2 L} \cdot \underbrace{\dfrac{\partial \widehat{y}_j}{\partial r_j^{(2)}}}_{\sigma'} \cdot \underbrace{\dfrac{\partial r_j^{(2)}}{\partial v_{ij}}}_{t_i}$$
于是隐藏层到输出层的权重就要做如下更新：

\pause

$$v_{ij} \quad{\color{red}\longrightarrow\quad} v_{ij} + h_{ij} = v_{ij} - \eta \cdot \partial_2 L \cdot \sigma' \cdot t_i$$


\pause

同理，输入层到隐藏层的权重要做如下更新：
$$w_{ij} \quad{\color{red}\longrightarrow\quad} w_{ij} - \eta \partial_2 L \cdot \sigma' \cdot x_i$$
\end{block}

\end{frame}

%------------------------------------------------

\begin{frame}

\begin{block}{反向传播}
\begin{itemize}
\item 通过这种反向的反馈机制，用输出的结果，反过来调节权重的大小。
\vspace{1em}
\pause
\item 用调整后的神经网络输入下一组样本，进行实验。用同样的方法再次调节权重的大小。依次重复，直到达到停止条件。
\end{itemize}

\end{block}

\end{frame}

% %------------------------------------------------

\begin{frame}

\begin{algorithm}[H]
\SetAlgoLined
\DontPrintSemicolon
\SetKwInOut{Input}{Input}
\SetKwInOut{Output}{Output}
\SetKwRepeat{Do}{do}{while}

\Input{训练集$\left\{\mathbf{y}_i, \mathbf{x}_i\right\}_{i=1}^N$}
\Input{可微凸损失函数$L(y,\widehat{y})$, 激励函数$\sigma$, 学习率$\eta$}
% \Input{weak learner $f_m(\mathbf{x}), m = 1,\ldots,M$}
% initialize model with a constant value
任选一组参数$(v_{ij}), (w_{ij}), b_i, b'_i$\;
$i \gets 1$ \;
\Do{达到停止条件}
{
输入$\mathbf{x}_i$, 计算预测值$\widehat{\mathbf{y}}_i$, 并保存中间结果（隐藏层$\mathbf{t}_i$）\;
计算$\partial_2 L, \sigma'$在相应自变量（$\widehat{\mathbf{y}}_i, \mathbf{t}_i$）处的取值 \;
更新隐藏层到输出层的权重以及Bias \;
更新输入层到隐藏层的权重以及Bias \;
$i \gets i+1$ \;
}
\Output{带参数$(v_{ij}), (w_{ij}), b_i, b'_i$的神经网络}
\caption{最简单的反向传播算法}
\end{algorithm}

\end{frame}

% %------------------------------------------------
% % closing page
% \begin{frame}
% \HUGE{\centerline{\bfseries {\color{gray} The} {\color{zkblue} End}}}

% \pause

% \vspace{1.2em}

% \Large{\centerline{\bfseries {\color{gray}Thank} \quad {\color{zkblue} You}}}

% \end{frame}

%----------------------------------------------------------------------------------------

\end{document}
